\documentclass[11pt]{article}

\usepackage[utf8]{inputenc}
\usepackage[margin=1in]{geometry}
\usepackage{amsmath,amssymb,mathtools,enumitem,fancyhdr}
\usepackage[misc]{ifsym}
\usepackage{hyperref}

\setlength\textwidth{7in}
\setlength\parindent{0pt}
\oddsidemargin=-0.5cm
\textheight=665pt
\pagestyle{fancy}
\setlength{\headheight}{13.6pt}
\renewcommand{\headrulewidth}{0.1pt}
\renewcommand{\footrulewidth}{0.1pt}
\lhead{\scshape Pontificia Universidad Católica de Chile}
\rhead{\scshape IMC UC}
\cfoot{\thepage}

\newcommand{\solution}{\large\textbf{Solución}\normalsize}
\newcommand{\norm}[1]{\left\lVert #1\right\rVert}
\newcommand{\abs}[1]{\left|#1\right|}
\def\mat   #1{\mbox{\boldmath $\mathsf #1$}}
\def\vec   #1{\mbox{\boldmath $#1$}{}}
\DeclareMathOperator{\grad}{grad}
\DeclareMathOperator{\dive}{div}
\def\R{\mathbb R}
\def\dx{\mbox{\(\,\mathrm{d}x\)}}
\def\ds{\mbox{\(\,\mathrm{d}s\)}}

\AtBeginDocument{\vspace*{-3.2\baselineskip}}

\begin{document}
    \begin{flushleft}
        \small
        \vspace{1pt}
        \textbf{IMT3130} \hfill
        \textbf{17/04/2026}
    \end{flushleft}
    \begin{center}
        \LARGE\textbf{Pauta Ayudantía 5}\\
        \vspace{3pt}
        \normalsize\textbf{Lifting, Galerkin y estimaciones abstractas}\\
        \normalsize Ayudante: Bastián Herrera \Letter{\hspace{1pt}: \texttt{biherrera@uc.cl}}
        \rule{\textwidth}{0.05mm}
    \end{center}

\section*{Problemas y soluciones}

\begin{enumerate}[label=\textbf{\arabic*.}, leftmargin=*]

    \item \textbf{Enunciado.} Considere
    \[
    \begin{aligned}
        -\Delta u+cu&=f &&\text{en }\Omega,\\
        \gamma_0u&=g &&\text{en }\Gamma_D,\\
        \partial_nu+\beta\gamma_0u&=r &&\text{en }\Gamma_R,
    \end{aligned}
    \]
    con
    \[
        V=\{v\in H^1(\Omega):\gamma_0v=0\text{ en }\Gamma_D\},
        \qquad
        V_g=\{v\in H^1(\Omega):\gamma_0v=g\text{ en }\Gamma_D\}.
    \]
    Resuelva primero el caso $g=0$ y luego reduzca el caso no homogéneo usando un lifting $G$ de la traza de Dirichlet.

    \textbf{Solución.}
    \begin{enumerate}[label=(\alph*)]
        \item Si $g=0$, el espacio de solución y de test es
        \[
            V=\{v\in H^1(\Omega):\gamma_0v=0\text{ en }\Gamma_D\}.
        \]
        Multiplicando por $v\in V$ e integrando por partes,
        \begin{align*}
            \int_\Omega fv\dx
            &=
            \int_\Omega (-\Delta u+cu)v\dx\\
            &=
            \int_\Omega \nabla u\cdot\nabla v\dx
            -\int_{\Gamma_R}\partial_nu\,\gamma_0v\ds
            +\int_\Omega cuv\dx.
        \end{align*}
        Como $\partial_nu=r-\beta\gamma_0u$ en $\Gamma_R$,
        \[
            a(u,v)=\ell(v)\qquad\forall v\in V,
        \]
        con
        \[
            a(w,v)=
            \int_\Omega \nabla w\cdot\nabla v\dx
            +\int_\Omega cwv\dx
            +\int_{\Gamma_R}\beta\,\gamma_0w\,\gamma_0v\ds,
        \]
        y
        \[
            \ell(v)=\int_\Omega fv\dx+\langle r,\gamma_0v\rangle_{\Gamma_R}.
        \]
        La continuidad de $a$ y $\ell$ sigue de Cauchy--Schwarz, de $c,\beta\in L^\infty$, de $f\in L^2$, de $r\in H^{-1/2}(\Gamma_R)$ y de la continuidad de la traza. Si $|\Gamma_D|>0$, Poincaré generalizado da
        \[
            \norm{v}_{H^1(\Omega)}\le C\norm{\nabla v}_{L^2(\Omega)}
            \qquad\forall v\in V.
        \]
        Como $c,\beta\ge0$,
        \[
            a(v,v)\ge \norm{\nabla v}_{L^2(\Omega)}^2\ge C\norm{v}_{H^1(\Omega)}^2.
        \]
        Por Lax--Milgram, existe un único $u\in V$.

        \item Para $g\in H^{1/2}(\Gamma_D)$, tomamos un lifting $G\in H^1(\Omega)$ tal que
        \[
            \gamma_0G=g\quad\text{en }\Gamma_D.
        \]
        Buscamos $u=u_0+G$, con $u_0\in V$. Al reemplazar en la formulación débil,
        \[
            a(u_0+G,v)=\ell(v)\qquad\forall v\in V.
        \]
        Por tanto, $u_0$ satisface
        \[
            a(u_0,v)=\widetilde\ell(v)\qquad\forall v\in V,
        \]
        donde
        \[
            \widetilde\ell(v)=\ell(v)-a(G,v).
        \]
        Es decir,
        \[
        \begin{aligned}
            \widetilde\ell(v)
            &=
            \int_\Omega fv\dx+\langle r,\gamma_0v\rangle_{\Gamma_R}
            -\int_\Omega \nabla G\cdot\nabla v\dx
            -\int_\Omega cGv\dx
            -\int_{\Gamma_R}\beta\,\gamma_0G\,\gamma_0v\ds.
        \end{aligned}
        \]
        Como $a$ es continua y coerciva en $V$, y $\widetilde\ell\in V'$, Lax--Milgram entrega un único $u_0\in V$. Entonces $u=u_0+G\in V_g$ existe y es único. En efecto, si dos soluciones en $V_g$ existieran, su diferencia pertenecería a $V$ y resolvería el problema homogéneo, luego sería cero por coercividad.
    \end{enumerate}

    \item \textbf{Enunciado.} Sea $V$ Hilbert, $a:V\times V\to\R$ simétrica, continua y coerciva, y sea $V_h\subset V$. Usando la ortogonalidad de Galerkin,
    \[
        a(u-u_h,v_h)=0\qquad\forall v_h\in V_h,
    \]
    pruebe que $\|v\|_a=a(v,v)^{1/2}$ es equivalente a la norma de $V$, demuestre la propiedad de mejor aproximación y escriba el sistema lineal.

    \textbf{Solución.}
    \begin{enumerate}[label=(\alph*)]
        \item Como $a$ es simétrica, continua y coerciva,
        \[
            \alpha\norm{v}_V^2\le a(v,v)\le M\norm{v}_V^2.
        \]
        Por tanto,
        \[
            \sqrt{\alpha}\norm{v}_V\le \norm{v}_a\le \sqrt M\norm{v}_V,
        \]
        y $\norm{\cdot}_a$ es una norma equivalente a la norma de $V$.

        \item Sea $e=u-u_h$ y sea $v_h\in V_h$. Como $v_h-u_h\in V_h$, por ortogonalidad de Galerkin,
        \[
            a(e,v_h-u_h)=0.
        \]
        Entonces
        \begin{align*}
            \norm{u-v_h}_a^2
            &=
            \norm{(u-u_h)+(u_h-v_h)}_a^2\\
            &=
            \norm{u-u_h}_a^2
            +\norm{u_h-v_h}_a^2
            +2a(u-u_h,u_h-v_h)\\
            &=
            \norm{u-u_h}_a^2+\norm{u_h-v_h}_a^2
            \ge \norm{u-u_h}_a^2.
        \end{align*}
        Como la igualdad se obtiene con $v_h=u_h$,
        \[
            \norm{u-u_h}_a=\inf_{v_h\in V_h}\norm{u-v_h}_a.
        \]

        \item Si $u_h=\sum_{j=1}^N U_j\varphi_j$, entonces
        \[
            \sum_{j=1}^N a(\varphi_j,\varphi_i)U_j=\ell(\varphi_i),
            \qquad i=1,\dots,N.
        \]
        En forma matricial,
        \[
            \mat A\vec U=\vec F,
            \qquad
            A_{ij}=a(\varphi_j,\varphi_i),
            \qquad
            F_i=\ell(\varphi_i).
        \]
    \end{enumerate}

    \item \textbf{Enunciado.} Sea
    \[
        \mat A_\varepsilon(x)=\mat I+\varepsilon\mat B(x),
        \qquad
        0\le\varepsilon<\frac1M,\qquad
        \|\mat B\|_{L^\infty(\Omega)}\le M,
    \]
    con $\mat B$ simétrica. Considere
    \[
    \begin{aligned}
        -\dive(\mat A_\varepsilon\grad u)+cu&=f &&\text{en }\Omega,\\
        \gamma_0u&=g &&\text{en }\Gamma_D,\\
        (\mat A_\varepsilon\grad u)\cdot\vec n+\beta\gamma_0u&=r &&\text{en }\Gamma_R.
    \end{aligned}
    \]
    Use un lifting para derivar la formulación débil, pruebe continuidad/coercividad, recupere la forma fuerte bajo regularidad y discuta el caso Neumann puro.

    \textbf{Solución.}
    \begin{enumerate}[label=(\alph*)]
        \item Tomamos $G\in H^1(\Omega)$ tal que $\gamma_0G=g$ en $\Gamma_D$ y escribimos $u=u_0+G$, con $u_0\in V$. La forma bilineal natural es
        \[
            a_\varepsilon(w,v)=
            \int_\Omega \mat A_\varepsilon\grad w\cdot\grad v\dx
            +\int_\Omega cwv\dx
            +\int_{\Gamma_R}\beta\,\gamma_0w\,\gamma_0v\ds.
        \]
        La formulación equivalente para $u_0$ es
        \[
            a_\varepsilon(u_0,v)=\widetilde\ell_\varepsilon(v)
            \qquad\forall v\in V,
        \]
        donde
        \[
        \begin{aligned}
            \widetilde\ell_\varepsilon(v)
            &=
            \int_\Omega fv\dx
            +\langle r,\gamma_0v\rangle_{\Gamma_R}
            -a_\varepsilon(G,v)\\
            &=
            \int_\Omega fv\dx
            +\langle r,\gamma_0v\rangle_{\Gamma_R}
            -\int_\Omega \mat A_\varepsilon\grad G\cdot\grad v\dx
            -\int_\Omega cGv\dx
            -\int_{\Gamma_R}\beta\,\gamma_0G\,\gamma_0v\ds.
        \end{aligned}
        \]

        \item Para continuidad, usando $\norm{\mat A_\varepsilon}_{L^\infty}\le 1+\varepsilon M$,
        \[
            \abs{a_\varepsilon(w,v)}
            \le
            C_\varepsilon\norm{w}_{H^1(\Omega)}\norm{v}_{H^1(\Omega)}.
        \]
        En efecto, los términos de reacción y borde se controlan como en el problema 1 mediante $L^\infty$ y la traza.

        Para coercividad, como $\mat B$ es simétrica,
        \[
            \xi^\top\mat A_\varepsilon(x)\xi
            =
            |\xi|^2+\varepsilon\,\xi^\top\mat B(x)\xi
            \ge
            (1-\varepsilon M)|\xi|^2
            \qquad\forall \xi\in\R^d.
        \]
        Como $0\le\varepsilon<1/M$, el número $\alpha_\varepsilon=1-\varepsilon M$ es positivo. Entonces, para $v\in V$,
        \[
            a_\varepsilon(v,v)
            \ge
            \alpha_\varepsilon\norm{\grad v}_{L^2(\Omega)}^2.
        \]
        Usando $|\Gamma_D|>0$ y Poincaré,
        \[
            a_\varepsilon(v,v)
            \ge
            C\alpha_\varepsilon\norm{v}_{H^1(\Omega)}^2.
        \]
        Por tanto, $a_\varepsilon$ es coerciva. Como $\widetilde\ell_\varepsilon\in V'$, Lax--Milgram entrega un único $u_0\in V$. Luego $u=u_0+G$ es la única solución en $V_g$.

        \item Supongamos que la solución es regular. La formulación para $u\in V_g$ dice
        \[
            a_\varepsilon(u,v)=
            \int_\Omega fv\dx+\langle r,\gamma_0v\rangle_{\Gamma_R}
            \qquad\forall v\in V.
        \]
        Si $\varphi\in C_c^\infty(\Omega)$, entonces $\varphi\in V$ y no hay términos de borde:
        \[
            \int_\Omega \mat A_\varepsilon\grad u\cdot\grad\varphi\dx
            +\int_\Omega cu\varphi\dx
            =
            \int_\Omega f\varphi\dx.
        \]
        Por definición distribucional,
        \[
            -\dive(\mat A_\varepsilon\grad u)+cu=f
            \qquad\text{en }\Omega.
        \]
        Ahora, para $v\in V$, usando integración por partes y la ecuación fuerte,
        \begin{align*}
            a_\varepsilon(u,v)
            &=
            \int_\Omega\left[-\dive(\mat A_\varepsilon\grad u)+cu\right]v\dx
            +\langle(\mat A_\varepsilon\grad u)\cdot\vec n,\gamma_0v\rangle_{\partial\Omega}
            +\int_{\Gamma_R}\beta\,\gamma_0u\,\gamma_0v\ds\\
            &=
            \int_\Omega fv\dx
            +\langle(\mat A_\varepsilon\grad u)\cdot\vec n+\beta\gamma_0u,\gamma_0v\rangle_{\Gamma_R}.
        \end{align*}
        Al comparar con el lado derecho de la formulación débil,
        \[
            \langle(\mat A_\varepsilon\grad u)\cdot\vec n+\beta\gamma_0u-r,\gamma_0v\rangle_{\Gamma_R}=0
            \qquad\forall v\in V.
        \]
        Como las trazas sobre $\Gamma_R$ son arbitrarias dentro del espacio de traza correspondiente,
        \[
            (\mat A_\varepsilon\grad u)\cdot\vec n+\beta\gamma_0u=r
            \qquad\text{en }\Gamma_R.
        \]
        La condición $\gamma_0u=g$ en $\Gamma_D$ ya estaba incorporada en la elección del espacio afín $V_g$.

        \item Si $\Gamma_D=\varnothing$, $c=0$ y $\beta=0$, entonces
        \[
            a_\varepsilon(v,v)=\int_\Omega \mat A_\varepsilon\grad v\cdot\grad v\dx.
        \]
        Las constantes tienen gradiente cero, por lo que pertenecen al núcleo de la forma. No hay unicidad en $H^1(\Omega)$ y se necesita una condición de compatibilidad:
        \[
            \int_\Omega f\dx+\langle r,1\rangle_{\partial\Omega}=0.
        \]
        La formulación bien puesta se obtiene fijando la constante, por ejemplo en
        \[
            H^1_\ast(\Omega)=\left\{v\in H^1(\Omega):\int_\Omega v\dx=0\right\},
        \]
        donde Poincaré--Wirtinger recupera coercividad.
    \end{enumerate}

\end{enumerate}

\end{document}
